\documentclass[11pt]{article}
\usepackage[margin=25mm]{geometry}
\usepackage{fontspec}
\setmainfont{Latin Modern Roman}
\usepackage{amsmath,amssymb,mathtools}
\usepackage{unicode-math}
\setmathfont{Latin Modern Math}
\usepackage{xurl,hyperref}
\hypersetup{colorlinks=true,urlcolor=blue}
\setcounter{secnumdepth}{0}
\setlength{\parindent}{0pt}
\setlength{\parskip}{5pt}
\setlength{\emergencystretch}{3em}
\newcommand{\sref}[2]{\hyperlink{source:#1:#2}{[S#2]}}
\newcommand{\cataloguescope}[1]{\paragraph{Recorded target.}#1}
\newcommand{\R}{\mathbb R}
\newcommand{\C}{\mathbb C}
\newcommand{\N}{\mathbb N}
\newcommand{\Z}{\mathbb Z}
\newcommand{\Q}{\mathbb Q}
\newcommand{\F}{\mathbb F}
\newcommand{\eps}{\varepsilon}
\begin{document}
% UnsolvedMath ID 20002553; source catalogue code AIM-PHYSICS-0061
\setcounter{section}{20002552}
\section{Quantitative collinear finiteness and maximal-dimensional simplex rigidity}\label{20002553}
{\small UnsolvedMath ID 20002553 · Mathematical Physics}
\subsection{Definitions and mathematical statement}

\paragraph{Shared notation.}$\mathbb N=\{1,2,\ldots\}$; the standard number fields have their usual meanings. All additional conventions are stated locally.


\paragraph{Definitions and notation.}In the Newtonian $n$-body problem, positions $q_i\in\mathbb R^d$ and masses $m_i>0$ obey $m_i\ddot q_i=\nabla_{q_i}U$, where $U(q)=\sum_{i<j}m_im_j/|q_i-q_j|$ and the gravitational constant is normalized to one. Here $T=\tfrac12\sum_i m_i|\dot q_i|^2$, $E=T-U$, and a collision has $q_i=q_j$. The center of mass can be fixed at zero; moment of inertia is $I=\sum_i m_i|q_i|^2$. A central configuration has $\nabla_{q_i}U=-\lambda m_iq_i$ in center-of-mass coordinates for a common scalar $\lambda$. After fixing $I=1$, it is a critical point of $U$ on the collision-free inertia ellipsoid. Congruent configurations and scales are factored out for a finiteness count. A relative equilibrium is a rigidly rotating solution with time-independent shape. 

\paragraph{Catalogue wording retained for traceability.}
For the Newtonian $n$-body problem with $n\ge4$ and every choice of positive masses $m_i>0$, prove that the number of central configurations is finite, or give a counterexample, counting configurations modulo the standard translation, scale and rotational equivalences.
\paragraph{Scope and source qualification.}The source asks for all positive masses, not only generic masses or collinear configurations. The spatial configuration class must be stated; finiteness of normalized labeled configurations without rotation quotient would not be the intended statement. The catalogue title/status was paired with a machine-generated research attempt; that attempt is not peer-reviewed evidence of settlement. The statement covers the original source questions rather than only the title’s subcase.
\subsection{Short English statement}
Determine whether every positive-mass n-body system with n≥4 has finitely many central configurations up to the usual symmetries.
\subsection{Research attempt}
\paragraph{Model, date and investigated scope.}
OpenAI GPT-6 Astra (configured model \texttt{gpt-6-astra}), ultra reasoning; generated 11 October 2026. Proved one collinear central shape per labeled order and the n!/2 reversal quotient. General noncollinear finiteness for every positive mass choice remains unproved.

\paragraph{Sources and literature check.}
Read the supplied complete definition, catalogue wording, scope qualifications and source list for this canonical ID. The original-source citations are inherited as supplied references; no fresh primary-source or novelty search was performed for this attempt. No current open-status claim is inferred from this calculation.

\paragraph{Collinear finiteness from strict convexity.}Fix positive labeled masses and one ordering \(x_1<\cdots<x_n\) on a line. Consider
\[W(x)=\sum_{i<j}\frac{m_im_j}{x_j-x_i}+\frac12\sum_im_ix_i^2.
\]
The chamber is convex. Along a nonzero vector \(v\), its Hessian quadratic form is
\[D^2W[v,v]=\sum_{i<j}\frac{2m_im_j(v_j-v_i)^2}{(x_j-x_i)^3}+\sum_im_iv_i^2>0.
\]
Thus W is strictly convex. It tends to infinity when a collision boundary is approached with bounded coordinates, or when the coordinate norm tends to infinity, so it attains a unique minimum in the chamber. At that point \(\partial_iU=-m_ix_i\). Summing derivatives gives \(\sum m_ix_i=0\), so it is a centered collinear central configuration.
\paragraph{Counting up to scale and reversal.}Every collinear central configuration has a positive multiplier \(\lambda\), because \(\lambda I=U>0\). Scaling coordinates by \(\lambda^{1/3}\) changes the multiplier to one, putting it at the unique minimum for its ordering. Thus each labeled ordering gives exactly one shape. There are \(n!\) orderings; identifying reversal by a spatial rotation gives \(n!/2\) labeled collinear classes for \(n\ge2\). Additional identifications of equal-mass labels can only reduce that count and are a separate convention.
\paragraph{Full target remaining.}Strict convexity relies on fixed one-dimensional order and the convexity of reciprocal positive gaps. Noncollinear chambers do not have that structure, and the potential Hessian has negative tangential directions. This reconstructs the known collinear subcase for all positive masses, not finiteness of all normalized central configurations for \(n\ge4\). The first missing step is excluding continuous families or accumulation of noncollinear critical shapes after rotational reduction, including exceptional mass choices.

\paragraph{Verification and adversarial audit.}
Checked strict Hessian positivity, coercive boundary behavior, center of mass, multiplier scaling lambda\textasciicircum{}(1/3), reversal count and equal-label convention. This was an LLM adversarial audit, not independent human verification; no theorem prover was used. Prompts A then B applied to this ID: Checked strict Hessian positivity, coercive boundary behavior, center of mass, multiplier scaling lambda\textasciicircum{}(1/3), reversal count and equal-label convention. Scope remains limited: Exclude continuous or accumulating families of noncollinear central configurations after normalization and rotational reduction, including nongeneric masses.

Coordinator checked strict convexity, coercivity and the multiplier rescaling in the collinear proof, including the labeled reversal count. This additional check was model-assisted, not independent human verification.

\paragraph{Tools, process and contributions.}
Python 3 and local shell used to read supplied files and serialize this manuscript; mathematics derived symbolically in the text, with no numerical experiment or formal proof assistant. Contribution prompt A was followed by prompt B. GPT-6 Astra with ultra reasoning drafted the research and performed its stated self-audit; any additional coordinator contribution is identified in the verification and audit. A GPT-6-based Codex coordinator assembled the manuscript and checked repository integration; its exact variant and reasoning configuration were not exposed. No human mathematical contribution, independent human proof check or formal proof certification is supplied. Complete derivations and the consulted source references are included here; private reasoning transcripts are not published.

\paragraph{Final status and first remaining gap.}
\textbf{UNRESOLVED}. Exclude continuous or accumulating families of noncollinear central configurations after normalization and rotational reduction, including nongeneric masses.

\section{Completion Estimate}
\noindent\textbf{COMPLETION: 0\%}
The calculation is an elementary benchmark or conditional reduction; it establishes no claimed new portion of the complete source target.

\subsection{Sources}
\begin{itemize}
\item[\hypertarget{source:20002553:1}{S1}] ULAM.AI and the UnsolvedMath Contributors, UnsolvedMath, record 20002553 (AIM-PHYSICS-0061). Catalogue curation under CC BY 4.0; definitions and English exposition added locally.
\url{https://huggingface.co/datasets/ulamai/UnsolvedMath}.
\item[\hypertarget{source:20002553:2}{S2}] American Institute of Mathematics, Variational Methods in Celestial Mechanics; notes by Kuo-Chang Chen, version 22 June 2003, source item 24.
\url{https://aimath.org/WWN/varcelest/varcelest.pdf}.
\end{itemize}
\label{end:20002553}
\end{document}
